\chapter{2010年广东卷（文）}

\section{单选题}

\begin{problem}
若集合 \(A = \{0, 1, 2, 3\}\)，\(B = \{1, 2, 4\}\)，则集合 \(A \cup B =\)（\quad）

\choices
    {\(\{0, 1, 2, 3, 4\}\)}
    {\(\{1, 2, 3, 4\}\)}
    {\(\{1, 2\}\)}
    {\(\{0\}\)}
\end{problem}

\begin{answer}
A
\end{answer}

\begin{solution}
并集包含两个集合中的全部元素，因此
\(A\cup B=\{0,1,2,3,4\}\)，故选 A.
\end{solution}

\begin{problem}
函数 \(f(x) = \lg(x - 1)\) 的定义域是（\quad）

\choices
    {\((2, +\infty)\)}
    {\((1, +\infty)\)}
    {\([1, +\infty)\)}
    {\([2, +\infty)\)}
\end{problem}

\begin{answer}
B
\end{answer}

\begin{solution}
对数的真数必须为正数，所以定义域满足 \(x-1>0\)，即 \(x>1\)，故选 B.
\end{solution}

\begin{problem}
若函数 \(f(x) = 3^x + 3^{-x}\) 与 \(g(x) = 3^x - 3^{-x}\) 的定义域均为 \(\R\)，则（\quad）

\choices
    {\(f(x)\) 与 \(g(x)\) 均为偶函数}
    {\(f(x)\) 为奇函数，\(g(x)\) 为偶函数}
    {\(f(x)\) 与 \(g(x)\) 均为奇函数}
    {\(f(x)\) 为偶函数，\(g(x)\) 为奇函数}
\end{problem}

\begin{answer}
D
\end{answer}

\begin{solution}
\(f(-x)=3^{-x}+3^x=f(x)\)，所以 \(f(x)\) 是偶函数；又
\(g(-x)=3^{-x}-3^x=-g(x)\)，所以 \(g(x)\) 是奇函数，故选 D.
\end{solution}

\begin{problem}
已知数列 \(\{a_n\}\) 为等比数列，\(S_n\) 是它的前 \(n\) 项和，若 \(a_2 \cdot a_3 = 2a_1\)，且 \(a_4\) 与 \(2a_7\) 的等差中项为 \(\frac{5}{4}\)，则 \(S_5 =\)（\quad）

\choices
    {\(35\)}
    {\(33\)}
    {\(31\)}
    {\(29\)}
\end{problem}

\begin{answer}
C
\end{answer}

\begin{solution}
设等比数列的首项为 \(a_1\)，公比为 \(q\). 由 \(a_2a_3=2a_1\) 及第二个条件可知 \(a_1\ne0\)，从而
\(a_1^2q^3=2a_1\)，即 \(a_1q^3=2\). 因此 \(a_4=a_1q^3=2\)，且
\(a_7=a_1q^6=(a_1q^3)q^3=2q^3=\dfrac{4}{a_1}\). 由等差中项的定义，
\(\dfrac{a_4+2a_7}{2}=\dfrac54\)，所以 \(\dfrac{2+8/a_1}{2}=\dfrac54\)，解得 \(a_1=16\). 于是 \(q^3=\dfrac18\)，即 \(q=\dfrac12\)，故
\(S_5=16\dfrac{1-(1/2)^5}{1-1/2}=31\)，故选 C.
\end{solution}

\begin{problem}
若向量 \(\bs{a} = (1, 1)\)，\(\bs{b} = (2, 5)\)，\(\bs{c} = (3, x)\) 满足条件 \(\left(8\bs{a} - \bs{b}\right) \cdot \bs{c} = 30\)，则 \(x =\)（\quad）

\choices
    {\(6\)}
    {\(5\)}
    {\(4\)}
    {\(3\)}
\end{problem}

\begin{answer}
C
\end{answer}

\begin{solution}
\(8\bs a-\bs b=(8,8)-(2,5)=(6,3)\). 由数量积条件，
\((8\bs a-\bs b)\cdot\bs c=(6,3)\cdot(3,x)=18+3x=30\)，所以 \(x=4\)，故选 C.
\end{solution}

\begin{problem}
若圆心在 \(x\) 轴上，半径为 \(\sqrt{5}\) 的圆 \(O\) 位于 \(y\) 轴左侧，且与直线 \(x + 2y = 0\) 相切，则圆 \(O\) 的方程是（\quad）

\choices
    {\(\left(x - \sqrt{5}\right)^2 + y^2 = 5\)}
    {\(\left(x + \sqrt{5}\right)^2 + y^2 = 5\)}
    {\((x - 5)^2 + y^2 = 5\)}
    {\((x + 5)^2 + y^2 = 5\)}
\end{problem}

\begin{answer}
D
\end{answer}

\begin{solution}
设圆心为 \(O(h,0)\). 因圆位于 \(y\) 轴左侧，\(h<0\). 圆心到切线 \(x+2y=0\) 的距离等于半径，故
\(\dfrac{|h|}{\sqrt5}=\sqrt5\)，从而 \(|h|=5\)，结合 \(h<0\) 得 \(h=-5\). 所以圆的方程为
\((x+5)^2+y^2=5\)，故选 D.
\end{solution}

\begin{problem}
若一个椭圆长轴的长度，短轴的长度和焦距成等差数列，则该椭圆的离心率是（\quad）

\choices
    {\(\frac{4}{5}\)}
    {\(\frac{3}{5}\)}
    {\(\frac{2}{5}\)}
    {\(\frac{1}{5}\)}
\end{problem}

\begin{answer}
B
\end{answer}

\begin{solution}
设椭圆的半长轴、半短轴和半焦距分别为 \(a\)，\(b\)，\(c\)，则长轴长度、短轴长度和焦距分别为 \(2a\)，\(2b\)，\(2c\)，且 \(a^2=b^2+c^2\). 由题意
\(2a+2c=4b\)，即 \(a+c=2b\). 令离心率 \(e=\dfrac ca\)，则 \(b=a\sqrt{1-e^2}\)，所以
\(1+e=2\sqrt{1-e^2}\). 因为 \(0<e<1\)，平方并整理得 \(5e^2+2e-3=0\)，故 \(e=\dfrac35\)，选 B.
\end{solution}

\begin{problem}
“\(x > 0\)”是“\(\sqrt[3]{x^2} > 0\)”的（\quad）

\choices
    {充分非必要条件}
    {必要非充分条件}
    {非充分非必要条件}
    {充要条件}
\end{problem}

\begin{answer}
A
\end{answer}

\begin{solution}
当 \(x>0\) 时，\(x^2>0\)，所以 \(\sqrt[3]{x^2}>0\). 反过来，\(\sqrt[3]{x^2}>0\) 只要求 \(x\ne0\)，例如 \(x=-1\) 时也成立，不能推出 \(x>0\). 因此“\(x>0\)”是“\(\sqrt[3]{x^2}>0\)”的充分非必要条件，故选 A.
\end{solution}

\begin{problem}
如图，\(\triangle ABC\) 为正三角形，\(AA' \parallel BB' \parallel CC'\)，\(CC' \perp\) 平面 \(ABC\)，且 \(3AA' = \frac{3}{2} BB' = CC' = AB\)，则多面体 \(ABC - A'B'C'\) 的正视图（也称主视图）是（\quad）

\choices
    {\(\choicebitmap[width=0.15\paperwidth]{img/2010/guangdong_liberal/q6_optA.png}\)}
    {\(\choicebitmap[width=0.15\paperwidth]{img/2010/guangdong_liberal/q6_optB.png}\)}
    {\(\choicebitmap[width=0.15\paperwidth]{img/2010/guangdong_liberal/q6_optC.png}\)}
    {\(\choicebitmap[width=0.15\paperwidth]{img/2010/guangdong_liberal/q6_optD.png}\)}
\end{problem}

\begin{answer}
D
\end{answer}

\begin{solution}
由题图，\(CC'\perp\) 平面 \(ABC\)，且 \(AA'\parallel BB'\parallel CC'\)，所以三条侧棱均垂直于底面. 设 \(AB=a\)，由
\(3AA'=\dfrac32BB'=CC'=a\) 得
\(AA'=\dfrac a3\)，\(BB'=\dfrac{2a}{3}\)，\(CC'=a\). 在正投影方向下，点 \(C\) 的水平投影位于 \(AB\) 的中点，故正视图中从左到右三个顶点的高度依次为 \(\dfrac a3\)，\(a\)，\(\dfrac{2a}{3}\)；\(A'B'\) 是可见的内部棱，\(C'C\) 为隐藏的竖棱. 这种轮廓和虚线位置与选项 D 相符，故选 D.
\end{solution}

\begin{problem}
在集合 \(\{a, b, c, d\}\) 上定义两种运算 \(\oplus\) 和 \(\otimes\) 如下：

\begin{center}
\begin{tabular}{c@{\qquad}c}
\begin{tabular}{c|cccc}
\(\oplus\) & \(a\) & \(b\) & \(c\) & \(d\) \\
\hline
\(a\) & \(a\) & \(b\) & \(c\) & \(d\) \\
\(b\) & \(b\) & \(b\) & \(b\) & \(b\) \\
\(c\) & \(c\) & \(b\) & \(c\) & \(b\) \\
\(d\) & \(d\) & \(b\) & \(b\) & \(d\)
\end{tabular}
&
\begin{tabular}{c|cccc}
\(\otimes\) & \(a\) & \(b\) & \(c\) & \(d\) \\
\hline
\(a\) & \(a\) & \(a\) & \(a\) & \(a\) \\
\(b\) & \(a\) & \(b\) & \(c\) & \(d\) \\
\(c\) & \(a\) & \(c\) & \(c\) & \(a\) \\
\(d\) & \(a\) & \(d\) & \(a\) & \(d\)
\end{tabular}
\end{tabular}
\end{center}

那么 \(d \otimes (a \oplus c) =\)（\quad）

\choices
    {\(a\)}
    {\(b\)}
    {\(c\)}
    {\(d\)}
\end{problem}

\begin{answer}
A
\end{answer}

\begin{solution}
由第一个运算表，\(a\oplus c=c\). 再由第二个运算表查得 \(d\otimes c=a\)，所以
\(d\otimes(a\oplus c)=a\)，故选 A.
\end{solution}

\section{填空题}

\begin{problem}
某城市缺水问题比较突出，为了制定节水管理办法，对全市居民某年的月均用水量进行了抽样调查，其中 4 位居民的月均用水量分别为 \(x_1\)，\(\cdots\)，\(x_4\)（单位：吨）. 根据如图所示的程序框图，若 \(x_1\)，\(x_2\)，\(x_3\)，\(x_4\) 分别为 1， 1.5， 1.5， 2，则输出的结果 \(s\) 为\fillinblank{}.

\bitmapfigure[max width=0.60\linewidth]{img/2010/guangdong_liberal/q11_fig1.png}
\end{problem}

\begin{answer}
\(1.5\)
\end{answer}

\begin{solution}
框图先令 \(s_1=0\)，从 \(i=1\) 开始，在 \(i\le4\) 时把 \(x_i\) 加入 \(s_1\)，再令 \(s=\dfrac1i s_1\). 因此逐步得到
\(s_1=1\)，\(s=1\)；\(s_1=2.5\)，\(s=1.25\)；\(s_1=4\)，\(s=\dfrac43\)；\(s_1=6\)，\(s=1.5\). 此时 \(i=5>4\)，输出 \(s=1.5\).
\end{solution}

\begin{problem}
某市居民 2005--2009 年家庭年平均收入 \(x\)（单位：万元）与年平均支出 \(Y\)（单位：万元）的统计资料如下表所示：

\begin{center}
\begin{tabular}{c|ccccc}
年份 & 2005 & 2006 & 2007 & 2008 & 2009 \\
\hline
收入 \(x\) & 11.5 & 12.1 & 13 & 13.3 & 15 \\
支出 \(Y\) & 6.8 & 8.8 & 9.8 & 10 & 12
\end{tabular}
\end{center}

根据统计资料，居民家庭年平均收入的中位数是\fillinblank{}，家庭年平均收入与年平均支出有\fillinblank{}（填“正”或“负”）线性相关关系.
\end{problem}

\begin{answer}
\(13\)，正
\end{answer}

\begin{solution}
收入数据已经按从小到大排列为 \(11.5\)，\(12.1\)，\(13\)，\(13.3\)，\(15\)，共有 5 个数据，所以中位数为中间的数 \(13\). 进一步，\(\bar x=12.98\)，\(\bar Y=9.48\)，并且
\(\sum_{i=1}^5(x_i-\bar x)(Y_i-\bar Y)=(-1.48)(-2.68)+(-0.88)(-0.68)+(0.02)(0.32)+(0.32)(0.52)+(2.02)(2.52)=9.828>0\). 协方差为正，说明收入与支出具有正线性相关关系.
\end{solution}

\begin{problem}
已知 \(a\)，\(b\)，\(c\) 分别是 \(\triangle ABC\) 的三个内角 \(A\)，\(B\)，\(C\) 所对的边. 若 \(a = 1\)，\(b = \sqrt{3}\)，\(A + C = 2B\)，则 \(\sin A =\)\fillinblank{}.
\end{problem}

\begin{answer}
\(\frac{1}{2}\)
\end{answer}

\begin{solution}
三角形内角和给出 \(A+B+C=\pi\). 结合 \(A+C=2B\)，得 \(3B=\pi\)，所以 \(B=\dfrac\pi3\). 由正弦定理，
\(\dfrac{a}{\sin A}=\dfrac{b}{\sin B}\)，故
\(\sin A=\dfrac{a}{b}\sin B=\dfrac1{\sqrt3}\cdot\dfrac{\sqrt3}{2}=\dfrac12\).
\end{solution}

\section{选考题：填空题}

\begin{problem}
如图，在直角梯形 \(ABCD\) 中，\(DC \parallel AB\)，\(CB \perp AB\)，\(AB = AD = a\)，\(CD = \frac{a}{2}\)，点 \(E\)，\(F\) 分别为线段 \(AB\)，\(AD\) 的中点，则 \(EF =\)\fillinblank{}.
\end{problem}

\begin{answer}
\(\frac{a}{2}\)
\end{answer}

\begin{solution}
取 \(A=(0,0)\)，\(B=(a,0)\). 因为 \(CB\perp AB\)，设 \(C=(a,h)\)；又因 \(DC\parallel AB\) 且 \(CD=\dfrac a2\)，按梯形顶点顺序可设 \(D=(\dfrac a2,h)\). 由 \(AD=a\) 得
\(\dfrac{a^2}{4}+h^2=a^2\)，所以 \(h=\dfrac{\sqrt3}{2}a\). 因而
\(E=(\dfrac a2,0)\)，\(F=(\dfrac a4,\dfrac{\sqrt3}{4}a)\)，从而
\(EF=\sqrt{\left(\dfrac a4\right)^2+\left(\dfrac{\sqrt3}{4}a\right)^2}=\dfrac a2\).
\end{solution}

\begin{problem}
在极坐标系 \((\rho, \theta)\)（\(0 \leqslant \theta < 2\pi\)）中，曲线 \(\rho(\cos \theta + \sin \theta) = 1\) 与 \(\rho(\sin \theta - \cos \theta) = 1\) 的交点的极坐标为\fillinblank{}.
\end{problem}

\begin{answer}
\(\left(1, \frac{\pi}{2}\right)\)
\end{answer}

\begin{solution}
两式相加、相减分别得到 \(2\rho\sin\theta=2\) 和 \(2\rho\cos\theta=0\). 因 \(\rho\sin\theta=1\)，有 \(\rho\ne0\)，所以 \(\cos\theta=0\). 在 \(0\leqslant\theta<2\pi\) 内，候选角为 \(\dfrac\pi2\) 和 \(\dfrac{3\pi}{2}\). 由 \(\rho\sin\theta=1\) 及极坐标中 \(\rho\ge0\)，只有 \(\theta=\dfrac\pi2\)，且 \(\rho=1\). 故交点的极坐标为 \(\left(1,\dfrac\pi2\right)\).
\end{solution}

\section{解答题}

\begin{problem}
设函数 \(f(x) = 3\sin \left(\omega x + \frac{\pi}{6}\right)\)（\(\omega > 0\)，\(x \in (-\infty, +\infty)\)），且以 \(\frac{\pi}{2}\) 为最小正周期.

\begin{enumerate}
\item 求 \(f(0)\)；
\item 求 \(f(x)\) 的解析式；
\item 已知 \(f\left(\frac{\alpha}{4} + \frac{\pi}{12}\right) = \frac{9}{5}\)，求 \(\sin \alpha\) 的值.
\end{enumerate}
\end{problem}

\begin{solution}
\begin{enumerate}
\item 正弦函数的最小正周期为 \(2\pi\)，所以 \(f(x)\) 的最小正周期为 \(\dfrac{2\pi}{\omega}\). 由题意 \(\dfrac{2\pi}{\omega}=\dfrac\pi2\)，得 \(\omega=4\)，从而 \(f(0)=3\sin\dfrac\pi6=\dfrac32\).
\item 由上面的 \(\omega=4\)，得 \(f(x)=3\sin\left(4x+\dfrac\pi6\right)\).
\item 代入已知条件，得
\(3\sin\left(4\left(\dfrac\alpha4+\dfrac\pi{12}\right)+\dfrac\pi6\right)=\dfrac95\)，即 \(3\sin\left(\alpha+\dfrac\pi2\right)=\dfrac95\). 因为 \(\sin(\alpha+\dfrac\pi2)=\cos\alpha\)，所以 \(\cos\alpha=\dfrac35\)，从而 \(\sin^2\alpha=1-\dfrac9{25}=\dfrac{16}{25}\)，故 \(\sin\alpha=\pm\dfrac45\).
\end{enumerate}
\end{solution}

\begin{problem}
某电视台在一次对收看文艺节目和新闻节目观众的抽样调查中，随机抽取了 100 名电视观众，相关的数据如下表所示：

\begin{center}
\begin{tabular}{c|ccc}
 & 文艺节目 & 新闻节目 & 总计 \\
\hline
\(20\) 到 \(40\) 岁 & 40 & 18 & 58 \\
大于 \(40\) 岁 & 15 & 27 & 42 \\
\hline
总计 & 55 & 45 & 100
\end{tabular}
\end{center}

\begin{enumerate}
\item 由表中数据直观分析，收看新闻节目的观众是否与年龄有关？
\item 用分层抽样方法在收看新闻节目的观众中随机抽取 5 名，大于 \(40\) 岁的观众应该抽取几名？
\item 在上述抽取的 5 名观众中任取 2 名，求恰有 1 名观众的年龄为 \(20\) 至 \(40\) 岁的概率.
\end{enumerate}
\end{problem}

\begin{solution}
\begin{enumerate}
\item 在 20 至 40 岁观众中收看新闻节目的比例为 \(\dfrac{18}{58}=\dfrac9{29}\)，在大于 40 岁观众中收看新闻节目的比例为 \(\dfrac{27}{42}=\dfrac9{14}\). 两个比例不同，所以收看新闻节目与年龄有关.
\item 收看新闻节目的观众共有 45 人，其中大于 40 岁的有 27 人. 按年龄分层抽取 5 人，应抽取 \(5\times\dfrac{27}{45}=3\) 名大于 40 岁的观众.
\item 上述 5 人中有 2 名年龄为 20 至 40 岁，有 3 名大于 40 岁. 任取 2 人，恰有 1 名年龄为 20 至 40 岁的概率为
\(\dfrac{\mathrm C_2^1\mathrm C_3^1}{\mathrm C_5^2}=\dfrac35\).
\end{enumerate}
\end{solution}

\begin{problem}
如图，\(\widehat{AEC}\) 是半径为 \(a\) 的半圆，\(AC\) 为直径，点 \(E\) 为 \(\widehat{AC}\) 的中点，点 \(B\) 和点 \(C\) 为线段 \(AD\) 的三等分点，平面 \(AEC\) 外一点 \(F\) 满足 \(FC \perp\) 平面 \(BED\)，\(FB = \sqrt{5}a\).

\begin{enumerate}
\item 证明：\(EB \perp FD\)；
\item 求点 \(B\) 到平面 \(FED\) 的距离.
\end{enumerate}

\FigureLayoutDeclare{img/2010/guangdong_liberal/q18_fig1.png}{10.6cm}{14bp 30bp 21bp 20bp}
\bitmapfigure[max width=0.45\linewidth]{img/2010/guangdong_liberal/q18_fig1.png}
\end{problem}

\begin{solution}
\begin{enumerate}
\item 取平面 \(AEC\) 为 \(z=0\)，令 \(A=(0,0,0)\)，\(C=(2a,0,0)\)，\(E=(a,a,0)\). 因 \(B\)，\(C\) 是线段 \(AD\) 的三等分点，且 \(AC=2a\)，有 \(AB=BC=CD=a\)，故 \(B=(a,0,0)\)，\(D=(3a,0,0)\). 由于 \(A\)，\(B\)，\(C\)，\(D\) 共线且 \(E\) 不在直线 \(AC\) 上，平面 \(BED\) 就是平面 \(AEC\). 又因为 \(FC\perp\) 平面 \(BED\)，可令 \(F=(2a,0,h)\). 由 \(FB=\sqrt5a\)，得 \(a^2+h^2=5a^2\)，所以 \(|h|=2a\)，取 \(h=2a\). 于是 \(\overrightarrow{EB}=(0,-a,0)\)，\(\overrightarrow{FD}=(a,0,-2a)\)，从而 \(\overrightarrow{EB}\cdot\overrightarrow{FD}=0\)，故 \(EB\perp FD\).
\item 由 \(\overrightarrow{ED}=(2a,-a,0)\)，\(\overrightarrow{EF}=(a,-a,2a)\)，可取平面 \(FED\) 的一个法向量为 \((2,4,1)\). 因此平面 \(FED\) 的方程为 \(2x+4y+z=6a\). 点 \(B=(a,0,0)\) 到该平面的距离为
\(\dfrac{|2a-6a|}{\sqrt{2^2+4^2+1^2}}=\dfrac{4a}{\sqrt{21}}=\dfrac{4\sqrt{21}}{21}a\).
\end{enumerate}
\end{solution}

\begin{problem}
某营养师要为某个儿童预定午餐和晚餐. 已知一个单位的午餐含 12 个单位的碳水化合物，6 个单位的蛋白质和 6 个单位的维生素 \(C\)；一个单位的晚餐含 8 个单位的碳水化合物，6 个单位的蛋白质和 10 个单位的维生素 \(C\). 另外，该儿童这两餐需要的营养中至少含 64 个单位的碳水化合物，42 个单位的蛋白质和 54 个单位的维生素 \(C\). 如果一个单位的午餐，晚餐的费用分别是 \(2.5\text{ 元}\) 和 \(4\text{ 元}\)，那么要满足上述的营养要求，并且花费最少，应当为该儿童分别预订多少个单位的午餐和晚餐？
\end{problem}

\begin{solution}
设午餐和晚餐分别预订 \(x\) 和 \(y\) 个单位，则 \(x\)，\(y\ge0\)，营养约束为
\[
\begin{cases}
12x+8y\ge64,\\
6x+6y\ge42,\\
6x+10y\ge54.
\end{cases}
\]
费用为 \(C=2.5x+4y\). 将约束分别化为 \(3x+2y\ge16\)，\(x+y\ge7\)，\(3x+5y\ge27\). 在 \(x=0\) 时，约束要求 \(y\ge8\)，得到边界顶点 \((0,8)\)；在 \(y=0\) 时，约束要求 \(x\ge9\)，得到边界顶点 \((9,0)\). 两条相邻边界的交点满足
\[
\begin{cases}
3x+2y=16,\\
x+y=7,
\end{cases}
\]
解得 \((x,y)=(2,5)\)，满足第三个约束；满足后两个等式的交点满足
\[
\begin{cases}
x+y=7,\\
3x+5y=27,
\end{cases}
\]
解得 \((x,y)=(4,3)\)，满足第一个约束. 第一个和第三个等式的交点为 \((26/9,11/3)\)，但此时 \(x+y=59/9<7\)，不在可行域内. 因而可行域的边界顶点为 \((0,8)\)，\((2,5)\)，\((4,3)\)，\((9,0)\). 在这些顶点处，费用依次为 \(32\)，\(25\)，\(22\)，\(22.5\). 因费用函数在可行域上的最小值必在顶点取得，故最少费用在 \((x,y)=(4,3)\) 处取得. 应分别预订 4 个单位的午餐和 3 个单位的晚餐.
\end{solution}

\begin{problem}
已知函数 \(f(x)\) 对任意实数 \(x\) 均有 \(f(x) = kf(x + 2)\)，其中常数 \(k\) 为负数，且 \(f(x)\) 在区间 \([0, 2]\) 上有表达式 \(f(x) = x(x - 2)\).

\begin{enumerate}
\item 求 \(f(-1)\)，\(f(2.5)\) 的值；
\item 写出 \(f(x)\) 在 \([-3, 3]\) 上的表达式，并讨论函数 \(f(x)\) 在 \([-3, 3]\) 上的单调性；
\item 求出 \(f(x)\) 在 \([-3, 3]\) 上的最小值与最大值，并求出相应的自变量的取值.
\end{enumerate}
\end{problem}

\begin{solution}
\begin{enumerate}
\item 由 \(f(x)=kf(x+2)\)，当 \(-2\le x\le0\) 时有 \(f(x)=kx(x+2)\)，故 \(f(-1)=-k\). 当 \(2\le x\le3\) 时，\(f(x)=\dfrac1k f(x-2)=\dfrac1k(x-2)(x-4)\)，故 \(f(2.5)=-\dfrac{3}{4k}\).
\item
\examdisplaycases{f(x)=}{
k^2(x+2)(x+4), & -3\leqslant x< -2,\\
kx(x+2), & -2\leqslant x<0,\\
x(x-2), & 0\leqslant x<2,\\
\dfrac{1}{k}(x-2)(x-4), & 2\leqslant x\leqslant 3,
}
其中各段分别由 \(f(x)=k f(x+2)\) 反复代入得到. 各段导数的符号为：在 \([-3,-2)\) 上为 \(2k^2(x+3)\ge0\)，在 \([-2,0)\) 上为 \(2k(x+1)\)，在 \([0,2)\) 上为 \(2(x-1)\)，在 \([2,3]\) 上为 \(\dfrac{2(x-3)}k\ge0\). 又 \(f(-2)=f(0)=f(2)=0\)，所以函数在 \([-3,-1]\) 与 \([1,3]\) 上为增函数，在 \([-1,1]\) 上为减函数.
\item
由单调性，极值只可能在上述单调区间的端点取得；分点的函数值为 \(f(-2)=f(0)=f(2)=0\)，且 \(f(-3)=-k^2<0\)，\(f(1)=-1<0\)，\(f(-1)=-k>0\)，\(f(3)=-\dfrac1k>0\)，所以 \(0\) 不会成为最小值或最大值，只需比较这四个非零端点值：

当 \(k<-1\) 时，\(-k^2<-1\)，且 \(-k>-\dfrac1k\)，所以最小值为 \(-k^2\)，在 \(x=-3\) 处取得；最大值为 \(-k\)，在 \(x=-1\) 处取得.

当 \(k=-1\) 时，最小值为 \(-1\)，在 \(x=-3\) 或 \(x=1\) 处取得；最大值为 \(1\)，在 \(x=-1\) 或 \(x=3\) 处取得.

当 \(-1<k<0\) 时，\(-1<-k^2\)，且 \(-\dfrac1k>-k\)，所以最小值为 \(-1\)，在 \(x=1\) 处取得；最大值为 \(-\dfrac1k\)，在 \(x=3\) 处取得.
\end{enumerate}
\end{solution}

\begin{problem}
已知曲线 \(C_n: y = nx^2\)，点 \(P_n(x_n, y_n)\)（\(x_n > 0\)，\(y_n > 0\)）是曲线 \(C_n\) 上的点（\(n = 1\)，\(2\)，\(\cdots\)）.

\begin{enumerate}
\item 试写出曲线 \(C_n\) 在点 \(P_n\) 处的切线 \(l_n\) 的方程，并求出 \(l_n\) 与 \(y\) 轴的交点 \(Q_n\) 的坐标；
\item 若原点 \(O(0, 0)\) 到 \(l_n\) 的距离与线段 \(P_nQ_n\) 的长度之比取得最大值，试求点 \(P_n\) 的坐标 \((x_n, y_n)\)；
\item 设 \(m\) 与 \(k\) 为两个给定的不同的正整数，\(x_n\) 与 \(y_n\) 是满足前一问条件的点 \(P_n\) 的坐标，证明：\(\sum_{n=1}^{s} \left| \sqrt{\frac{(m + 1)x_n}{2}} - \sqrt{(k + 1)y_n} \right| < \left| \sqrt{ms} - \sqrt{ks} \right|\)（\(s = 1\)，\(2\)，\(\cdots\)）.
\end{enumerate}
\end{problem}

\begin{solution}
\begin{enumerate}
\item 由 \(y=nx^2\) 得 \(y'=2nx\)，所以在 \(P_n(x_n,y_n)\) 处的切线为
\(l_n:y-y_n=2nx_n(x-x_n)\). 因 \(y_n=nx_n^2\)，令 \(x=0\) 得 \(y=-nx_n^2\)，故 \(Q_n=(0,-nx_n^2)\).
\item 将切线写成 \(2nx_nx-y-nx_n^2=0\)，则原点到切线的距离为
\(d_n=\dfrac{nx_n^2}{\sqrt{1+4n^2x_n^2}}\). 又
\(P_nQ_n=\sqrt{x_n^2+(2nx_n^2)^2}=x_n\sqrt{1+4n^2x_n^2}\). 因而二者之比为
\(R_n=\dfrac{nx_n}{1+4n^2x_n^2}\). 令 \(u=2nx_n>0\)，则 \(R_n=\dfrac{u}{2(1+u^2)}\). 由 \((u-1)^2\ge0\) 得 \(1+u^2\ge2u\)，所以 \(R_n\le\dfrac14\)，等号当且仅当 \(u=1\). 因此 \(x_n=\dfrac1{2n}\)，再由 \(y_n=nx_n^2\) 得 \(y_n=\dfrac1{4n}\)，即 \(P_n\left(\dfrac1{2n},\dfrac1{4n}\right)\).
\item 由上一问，
\[
\left|\sqrt{\frac{(m+1)x_n}{2}}-\sqrt{(k+1)y_n}\right|=\frac{|\sqrt{m+1}-\sqrt{k+1}|}{2\sqrt n}
\]
又对每个正整数 \(n\)，有
\[
2(\sqrt n-\sqrt{n-1})=\frac{2}{\sqrt n+\sqrt{n-1}}>\frac1{\sqrt n}
\]
其中最后一个不等号等价于 \(\sqrt n>\sqrt{n-1}\). 求和得
\(\sum_{n=1}^s\dfrac1{\sqrt n}<2\sqrt s\). 另外，因 \(m\ne k\)，不妨设 \(m>k\)；有理化可得
\[
\sqrt{m+1}-\sqrt{k+1}=\frac{m-k}{\sqrt{m+1}+\sqrt{k+1}}<\frac{m-k}{\sqrt m+\sqrt k}=\sqrt m-\sqrt k
\]
若 \(m<k\) 则交换二者，故
\(|\sqrt{m+1}-\sqrt{k+1}|<|\sqrt m-\sqrt k|\). 因此
\[
\sum_{n=1}^{s}\left|\sqrt{\frac{(m+1)x_n}{2}}-\sqrt{(k+1)y_n}\right|
<|\sqrt{m+1}-\sqrt{k+1}|\sqrt s
<|\sqrt m-\sqrt k|\sqrt s
=|\sqrt{ms}-\sqrt{ks}|
\]
所证不等式成立.
\end{enumerate}
\end{solution}
